\section{Relações de equivalência}\label{sec:relacoesDeEquivalencia}
\index{relação!de equivalência}
\index{partição}

Nesta seção, introduziremos as noções básicas de relação de equivalência, classe de equivalência, conjunto quociente e partição, bem como alguns resultados estruturais associados.

Historicamente, a noção abstrata de relação de equivalência é posterior a muitos de seus exemplos.
Um exemplo de grande influência é a congruência módulo $m$, já sistematizada por Gauss nas \emph{Disquisitiones Arithmeticae}: em linguagem atual, ela identifica números que têm o mesmo resto módulo $m$ e os organiza em classes de resíduos \cite{gauss1801disquisitiones}.

\begin{definition}[$\BST$](Relação de equivalência)\label{def:relacaoDeEquivalencia}\index{relação!de equivalência}
    Uma \emph{relação de equivalência} é um par $(A, \sim)$ em que $A$ é um conjunto e $\sim$ é uma relação binária em $A$ tal que para todo $x, y, z \in A$, valem:
    \begin{itemize}
        \item $x \sim x$ (reflexividade)
        \item Se $x \sim y$, então $y \sim x$ (simetria)
        \item Se $x \sim y$ e $y \sim z$, então $x \sim z$ (transitividade).
    \end{itemize}
\end{definition}

Todo conjunto admite uma relação de equivalência trivial, a saber, a relação de igualdade.

\begin{example}[$\BST$]\label{exa:relacaoDeIgualdade}
    Seja $A$ um conjunto.
    A \emph{relação de igualdade em $A$} é a relação binária $=_A=\{(x, y) \in A \times A : x=y\}$.
    Então $(A, =_A)$ é uma relação de equivalência.
    Denotamos $(A, =_A)$ como $(A, =)$.
\end{example}

A relação de equivalência acima é mínima: se $R$ é uma relação de equivalência em $A$, então $=_A \subseteq R$.
Por outro lado, há a relação de equivalência máxima, a saber, a relação universal.

\begin{example}[$\BST$]\label{exa:relacaoUniversal}
    Seja $A$ um conjunto.
    Então $(A, A \times A)$ é uma relação de equivalência denominada \emph{relação universal em $A$}.
\end{example}

Toda função induz naturalmente uma relação de equivalência.
\begin{example}[$\BST$]\label{exa:relacaoDeEquivalenciaInduzidaPorFuncao}
    Sejam $A$ e $B$ conjuntos e seja $f : A \to B$ uma função.
    Defina a relação binária $\sim_f$ em $A$ por
    \[
    x \sim_f y \iff f(x) = f(y).
    \]
    Então $(A, \sim_f)$ é uma relação de equivalência.
\end{example}

Dada uma relação de equivalência, é natural agrupar os elementos equivalentes entre si.

\begin{definition}[$\BST$](Classe de equivalência)\label{def:classeDeEquivalencia}\index{classe de equivalência}\index{conjunto quociente}
    Seja $(A, \sim)$ uma relação de equivalência e seja $x \in A$.
    A \emph{classe de equivalência} de $x$ é o conjunto
    \[
    [x]_{\sim} = \{y \in A : y \sim x\}.
    \]

    ($Z^{-}$) O conjunto de todas as classes de equivalência de elementos de $A$ é denotado por $A/{\sim}$.
\end{definition}

Deixamos os seguintes fatos básicos para verificação do leitor (Exercício~\ref{exr:propriedadesBasicasDasClassesDeEquivalencia}).

\begin{proposition}[$\BST$]\label{prop:propriedadesBasicasDasClassesDeEquivalencia}
    Seja $(A, \sim)$ uma relação de equivalência e sejam $x, y \in A$. Então:
    \begin{enumerate}[label=(\alph*)]
        \item $x \in [x]_{\sim}$;
        \item $x \sim y$ se, e somente se, $y \in [x]_{\sim}$;
        \item $x \sim y$ se, e somente se, $[x]_{\sim}=[y]_{\sim}$;
        \item se $[x]_{\sim} \cap [y]_{\sim} \neq \emptyset$, então $[x]_{\sim}=[y]_{\sim}$;
        \item ou $[x]_{\sim}=[y]_{\sim}$, ou $[x]_{\sim} \cap [y]_{\sim}=\emptyset$.
        \item ($Z^{-}$) $\displaystyle A=\bigcup \Big(A/{\sim}\Big)$.
    \end{enumerate}
\end{proposition}

A Proposição~\ref{prop:propriedadesBasicasDasClassesDeEquivalencia} mostra que duas classes de equivalência diferentes são disjuntas, e que, reunidas, elas formam o conjunto $A$.
Assim, podemos imaginar que uma relação de equivalência \emph{particiona} o conjunto $A$ em blocos disjuntos.
Convém definir o conceito de partição para tornar essa noção mais precisa.

\begin{definition}[$\BST$](Partição)\label{def:particao}\index{partição}
    Seja $A$ um conjunto.
    Uma \emph{partição} de $A$ é um conjunto $\mathcal{E}$ tal que:
    \begin{enumerate}[label=(\roman*)]
        \item se $X \in \mathcal{E}$, então $X \neq \emptyset$ e $X \subseteq A$;
        \item $\bigcup \mathcal{E}=A$;
        \item se $X, Y \in \mathcal{E}$ e $X \cap Y \neq \emptyset$, então $X=Y$.
    \end{enumerate}
\end{definition}

Decorre da Proposição~\ref{prop:propriedadesBasicasDasClassesDeEquivalencia} que, dada uma relação de equivalência, o conjunto $A/{\sim}$ é uma partição de $A$.
Reciprocamente, temos a seguinte proposição.

\begin{proposition}[$Z^{-}$]\label{prop:particaoInduzRelacaoDeEquivalencia}
    Seja $A$ um conjunto e seja $\mathcal{E}$ uma partição de $A$.
    Defina a relação binária $\sim_{\mathcal{E}}$ em $A$ por
    \[
    x \sim_{\mathcal{E}} y \iff \exists X \in \mathcal{E}\, (x \in X \wedge y \in X).
    \]
    Então $(A, \sim_{\mathcal{E}})$ é uma relação de equivalência.
    Além disso, para todo $x \in A$, a classe de equivalência de $x$ é exatamente o único elemento de $\mathcal{E}$ que contém $x$.
    Em particular,
    \[
    A/{\sim_{\mathcal{E}}} = \mathcal{E}.
    \]
\end{proposition}

\begin{proof}
    Dado $x \in A$, seja $X \in \mathcal{E}$ tal que $x \in X$.
    Tal $X$ é único: se $x \in Y \in \mathcal{E}$, então $x \in X \cap Y$, e portanto $X=Y$.

    Para todo $x \in A$ existe $X \in \mathcal{E}$ tal que $x \in X$.
    Daí segue que $x \sim_{\mathcal{E}} x$.
    Logo, $\sim_{\mathcal{E}}$ é reflexiva.

    A simetria é imediata pela definição.

    Para a transitividade, suponha que $x \sim_{\mathcal{E}} y$ e $y \sim_{\mathcal{E}} z$.
    Então existem $X, Y \in \mathcal{E}$ tais que $x, y \in X$ e $y, z \in Y$.
    Pela unicidade do elemento de $\mathcal{E}$ que contém $y$, segue que $X=Y$, e, portanto, $x, z \in X$.
    Logo, $x \sim_{\mathcal{E}} z$.

    Agora fixe $x \in A$ e seja $X \in \mathcal{E}$ tal que $x \in X$.
    Mostremos que $[x]_{\sim_{\mathcal{E}}}=X$.
    Se $y \in [x]_{\sim_{\mathcal{E}}}$, então $y \sim_{\mathcal{E}} x$, logo existe $Y \in \mathcal{E}$ tal que $x, y \in Y$.
    Pela unicidade do elemento de $\mathcal{E}$ que contém $x$, segue que $Y=X$, e portanto $y \in X$.
    Assim, $[x]_{\sim_{\mathcal{E}}} \subseteq X$.
    Reciprocamente, se $y \in X$, então $x, y \in X$, e portanto $y \sim_{\mathcal{E}} x$.
    Logo, $y \in [x]_{\sim_{\mathcal{E}}}$.
    Isso conclui que $X=[x]_{\sim_{\mathcal{E}}}$.

    Portanto, os elementos de $A/{\sim_{\mathcal{E}}}$ são exatamente os elementos de $\mathcal{E}$, isto é, $A/{\sim_{\mathcal{E}}}=\mathcal{E}$.
\end{proof}

Dessa forma, toda relação de equivalência induz uma partição, e toda partição induz uma relação de equivalência.
Da Proposição~\ref{prop:particaoInduzRelacaoDeEquivalencia}, iniciando-se com uma relação de equivalência, passando-se para a partição induzida por ela e, depois, passando-se para a relação de equivalência induzida por essa partição, chega-se à relação de equivalência inicial.
A recíproca é dada a seguir.

\begin{proposition}[$Z^{-}$]\label{prop:relacaoDeEquivalenciaInduzidaPelaParticao}
    Seja $(A, \sim)$ uma relação de equivalência e seja $\mathcal{E}=A/{\sim}$ a partição de $A$ induzida por $\sim$.
    Então $\sim = \sim_{\mathcal{E}}$.
\end{proposition}
\begin{proof}
    Dados $x, y \in A$, temos
    \[
    \begin{aligned}
        x \sim_{\mathcal{E}} y
        &\iff \exists X \in A/{\sim}\, (x \in X \wedge y \in X) \\
        &\iff \exists z \in A\, (x \in [z]_{\sim} \wedge y \in [z]_{\sim}) \\
        &\iff x \sim y.
    \end{aligned}
    \]
\end{proof}

Assim, partições e relações de equivalência são duas maneiras de descrever o mesmo fenômeno matemático.

\subsection*{Exercícios}
\begin{exercise}\label{exr:propriedadesBasicasDasClassesDeEquivalencia}
    Prove a Proposição~\ref{prop:propriedadesBasicasDasClassesDeEquivalencia} (nas teorias indicadas).
\end{exercise}

\begin{exercise}[$\BST$]
    Seja $A$ um conjunto.
    Prove que $(A, =)$ e $(A, A \times A)$ são relações de equivalência.
    Determine, em cada caso, as classes de equivalência correspondentes.
\end{exercise}

\begin{exercise}[$\BST$]
    Seja $(A, \sim)$ uma relação de equivalência e seja $B \subseteq A$.
    Defina
    \[
    \sim_B = \sim \cap (B \times B).
    \]
    Prove que $(B, \sim_B)$ é uma relação de equivalência.
    Mostre também que, para todo $x \in B$, a classe de equivalência de $x$ em $(B, \sim_B)$ é $[x]_{\sim} \cap B$.
\end{exercise}

\begin{exercise}[$Z^-$]
    Seja $A$ um conjunto e seja $\sim$ uma relação de equivalência em $A$.
    Mostre que existe uma função $f$ com $\dom f = A$ tal que $\sim = \sim_f$.
\end{exercise}
